% Textbook-style overview of the two Kosović–Juliano 3D PBL closure variants
% Build: source ~/miniconda3/bin/activate tex && tectonic branko_juliano_schemes.tex
\documentclass[10pt,a4paper]{article}
\usepackage[hmargin=1.8cm,vmargin=1.6cm]{geometry}
\usepackage{amsmath,amssymb}
\usepackage{booktabs}
\usepackage{microtype}
\usepackage{titlesec}
\titlespacing*{\section}{0pt}{7pt plus 2pt minus 2pt}{4pt plus 1pt}
\titlespacing*{\subsection}{0pt}{5pt plus 2pt minus 2pt}{3pt plus 1pt}
\usepackage[hidelinks]{hyperref}
\usepackage{multicol}
\usepackage{xcolor}
\usepackage{array,tabularx}
% --- key assumptions for the PhD concept (Section \ref{sec:concept}) -------
\setlength{\marginparwidth}{1.30cm}
\setlength{\marginparsep}{5pt}
\definecolor{keyc}{HTML}{7B2D8E}
\newcommand{\keytag}[1]{\colorbox{keyc}{\textcolor{white}{\sffamily\bfseries\scriptsize K#1}}}
\newcommand{\key}[1]{\marginpar[\raggedleft\keytag{#1}]{\raggedright\keytag{#1}}\ignorespaces}
\newcommand{\kh}[1]{{\color{keyc!80!black}\bfseries #1}}

\AtBeginDocument{%
  \setlength{\abovedisplayskip}{4pt plus 1pt minus 2pt}%
  \setlength{\belowdisplayskip}{4pt plus 1pt minus 2pt}%
  \setlength{\abovedisplayshortskip}{2pt plus 1pt}%
  \setlength{\belowdisplayshortskip}{2pt plus 1pt}}

\newcommand{\avg}[1]{\langle #1 \rangle}
\newcommand{\pd}[2]{\frac{\partial #1}{\partial #2}}

\title{\vspace{-1.2em}The Kosović--Juliano Three-Dimensional PBL Scheme:\\
the Approximate and the Full Closure\vspace{-0.4em}}
\author{Notes for the Inn Valley campaign}
\date{2026-09-01; revised 2026-09-12 --- key assumptions of the PhD concept marked \keytag{n} in the margin, Section~\ref{sec:concept} added, surface-boundary paragraph corrected}

\begin{document}
\maketitle
\vspace{-2.2em}

\section{The problem both variants solve}

A conventional PBL parameterization closes the Reynolds-averaged equations under the assumption that each grid cell is horizontally homogeneous. Write the mean momentum and potential-temperature equations with the turbulent terms kept general,
\begin{equation}
\pd{U_i}{t} + U_j \pd{U_i}{x_j}
   = -\varepsilon_{ijk} f_j U_k - \frac{1}{\rho}\pd{P}{x_i}
     - \pd{\avg{u_i u_j}}{x_j} - g_i ,
\qquad
\pd{\Theta}{t} + U_j \pd{\Theta}{x_j} = -\pd{\avg{u_j \theta}}{x_j},
\label{eq:mean}
\end{equation}
where $U_i$ and $\Theta$ are the grid-cell-mean velocity (m\,s$^{-1}$) and potential temperature (K), $u_i$ and $\theta$ their subgrid fluctuations, $f_j$ the Coriolis vector (s$^{-1}$), $P$ the mean kinematic pressure (m$^2\,$s$^{-2}$), and angle brackets denote the grid-cell average; an equation of the form of the second one is carried for water vapour as well. The symmetric stress tensor $\avg{u_i u_j}$ (m$^2\,$s$^{-2}$) has six independent components and each scalar-flux vector $\avg{u_j\theta}$ (K\,m\,s$^{-1}$) three. \key{1}\kh{Under horizontal homogeneity all horizontal derivatives of these statistics vanish identically}, so of the nine stress-and-heat-flux quantities only three survive into the mean equations — $\avg{uw}$, $\avg{vw}$ and $\avg{w\theta}$ — \kh{and only through their \emph{vertical} derivative}. That is the entire content of a 1D PBL scheme: however sophisticated its closure, it delivers three numbers per level.

At $\Delta x = 500$~m over Alpine terrain the homogeneity assumption fails twice over. First, the grid spacing is comparable to the boundary-layer depth $z_i$ — the ``grey zone'' or \emph{terra incognita}, where the most energetic eddies are neither safely subgrid (the mesoscale limit $\Delta x \gg z_i$) nor resolved (the LES limit $\Delta x \ll z_i$). Horizontal gradients of the turbulence statistics are then of the same order as the vertical ones, and a model that ignores them while partially resolving convective structures produces the well-documented spurious circulations of the grey zone. Second, slopes of $30^\circ$ and more tilt the mean gradients so that ``horizontal'' and ``vertical'' cease to be distinguished directions of the flow: a katabatic jet generates shear \emph{across} the slope that projects onto every component of the strain tensor, and a closure that can only respond to $\partial U/\partial z$ and $\partial V/\partial z$ is structurally blind to most of the production acting on it.

The 3D PBL scheme of Kosović et al.\ (2020) and Juliano et al.\ (2022) responds by parameterizing \emph{all} six stress components and all three components of each scalar flux, and applying their full three-dimensional divergence to \eqref{eq:mean}. In the WRF implementation this is a \emph{dynamics}-level option (\texttt{pbl3d\_opt}, requiring \texttt{bl\_pbl\_physics}=0 and \texttt{diff\_opt}=0): \key{2}the scheme computes its own metric-corrected horizontal derivatives in the terrain-following coordinate and \kh{replaces both the 1D PBL vertical mixing and the Smagorinsky horizontal diffusion}. There is no eddy-viscosity backstop behind it — what the closure supplies is all the subgrid mixing the model has. The MYNN control keeps WRF's Smagorinsky operator, so the step from the control to either 3D variant changes the horizontal exchange as well as the closure.

The two variants — \texttt{pbl3d\_opt}$=\pm1$, the \emph{approximate} closure, and \texttt{pbl3d\_opt}$=2$, the \emph{full} closure — share every part of this machinery except one thing: \emph{how the second moments are obtained from the mean gradients}. They are two branches of the same module, not two schemes, which is why almost everything in Sections~\ref{sec:core}--\ref{sec:lscale} and \ref{sec:numerics} applies to both.

\section{The common engine: Mellor--Yamada level 2.5}
\label{sec:core}

\subsection{Where level 2.5 sits in the hierarchy}

Mellor and Yamada (1974) organized second-order closures into levels by how many moments remain prognostic. Level 4 carries prognostic equations for all second moments — ten transport equations plus the scalar variances, unwieldy in an NWP model. Level 3 keeps prognostic equations only for the turbulence kinetic energy and the temperature variance, making the remaining moments algebraic. Level 2.5, the workhorse of the family (MYNN belongs to it too), keeps a single prognostic variable, and level 2 keeps none: every moment, including the TKE, is diagnosed from a local equilibrium. The reduction at each step is the same pair of modelling ingredients applied more aggressively: Rotta's return-to-isotropy hypothesis for the pressure--strain correlations, and Kolmogorov's local isotropy for the dissipative terms, after which the material derivatives and transport of the demoted moments are neglected so that their equations become algebraic.

Both variants of the 3D scheme are level-2.5 truncations. The prognostic variable is
\begin{equation}
q^2 = \avg{u^2}+\avg{v^2}+\avg{w^2} = 2\,\mathrm{TKE}
\qquad (\text{m}^2\,\text{s}^{-2}),
\end{equation}
transported by
\begin{equation}
\pd{q^2}{t} + U_k \pd{q^2}{x_k}
 = \pd{}{x_k}\!\left( l\, q\, S_q \pd{q^2}{x_k} \right)
   + 2\left( P_s + P_b - \epsilon \right),
\label{eq:qsq}
\end{equation}
where the terms are, in order: advection by the resolved flow (handled by WRF's scalar transport, so $q^2$ is carried through the Runge--Kutta steps like a moist scalar); down-gradient turbulent self-transport with the diffusion-coefficient ratio $S_q = 0.2$ (Mellor--Yamada 1982, Eq.~24); shear production
\begin{equation}
P_s = -\avg{u_i u_j}\,\pd{U_i}{x_j};
\label{eq:ps}
\end{equation}
buoyancy production
\begin{equation}
P_b = \beta g \avg{w\theta_v},
\qquad \beta = \frac{1}{T_\mathrm{ref}},\quad T_\mathrm{ref}=300\ \text{K fixed},
\end{equation}
with $\theta_v$ the virtual potential temperature, so that an upward heat flux feeds the vertical velocity variance; and the dissipation rate
\begin{equation}
\epsilon = \frac{q^3}{B_1\, l}
\qquad (\text{m}^2\,\text{s}^{-3}).
\label{eq:eps}
\end{equation}
\key{3}Equation \eqref{eq:eps} is the signature of the family: \kh{$\epsilon$ is not modelled independently but diagnosed from the \emph{same} master length scale $l$ that sets the fluxes.} Dimensionally it is the Kolmogorov cascade estimate $u^3/\ell$ evaluated with the closure's own velocity and length scales; structurally it means one scale controls both how far an eddy transports and how fast it decays. Keeping that scale single-valued is not a stylistic preference — Section~\ref{sec:full} shows what happens when a limiter breaks it. (A level-2 option, \texttt{pbl3d\_prog}=0, replaces \eqref{eq:qsq} by the local balance $P_s + P_b = \epsilon$ solved for $q$; the default is prognostic.)

\subsection{The algebraic second-moment system}

\key{5}All remaining second moments are algebraic. \kh{Neglecting their material derivatives and transport} leaves the coupled system (Juliano et al.\ 2022, Eqs.~10--12)
\begin{align}
\avg{u_i u_j} &= \frac{\delta_{ij}}{3}q^2
 - \frac{3 l_1}{q}\Big[ \big(\avg{u_k u_i} - C_1 q^2 \delta_{ki}\big)\pd{U_j}{x_k}
 + \big(\avg{u_k u_j} - C_1 q^2 \delta_{kj}\big)\pd{U_i}{x_k}
 - \tfrac{2}{3}\delta_{ij} \avg{u_k u_n}\pd{U_n}{x_k} \Big] \notag \\
 &\quad - \frac{3 l_1}{q}\,\beta\Big[ g_j \avg{u_i \theta_v} + g_i \avg{u_j \theta_v}
 - \tfrac{2}{3}\delta_{ij}\, g_n \avg{u_n \theta_v} \Big],
\label{eq:stress} \\[3pt]
\avg{u_j \theta_v} &= -\frac{3 l_2}{q}\Big[ \avg{u_j u_k}\pd{\Theta_v}{x_k}
 + \avg{u_k \theta_v}\pd{U_j}{x_k} + \beta g_j \avg{\theta_v^2} \Big],
\label{eq:hflux} \\[3pt]
\avg{\theta_v^2} &= -\frac{\Lambda_2}{q}\avg{u_k \theta_v}\pd{\Theta_v}{x_k},
\label{eq:tvar}
\end{align}
and an analogous set for the water-vapour fluxes $\avg{u_j q_v}$ and the covariance $\avg{q_v\theta_v}$. The four length scales are proportional to one master scale,
\begin{equation}
(l_1,\ \Lambda_1,\ l_2,\ \Lambda_2) = l\,(A_1,\ B_1,\ A_2,\ B_2),
\end{equation}
with the closure constants of Mellor and Yamada (1982) as the default set (\texttt{pbl3d\_constants}=\texttt{'MY82'}):
\begin{equation}
(A_1,\ B_1,\ A_2,\ B_2,\ C_1) = (0.92,\ 16.6,\ 0.74,\ 10.1,\ 0.08).
\end{equation}
\key{7}\kh{These were calibrated against neutral surface-layer data}; the alternative sets M73, MY74, MYJ and the Bougeault--Lacarrère-tuned values $(0.3, 8.4, 0.33, 6.4, 0.08)$ are selectable (Juliano et al.\ published with the Bougeault--Lacarrère-tuned set; this project runs MY82). Note they differ throughout from MYNN's Nakanishi--Niino (2009) set — \kh{$B_1 = 24$ there against $16.6$ here}, so the two schemes do not even dissipate at the same rate for the same $q$ and $l$: at equal $q$ and $l$ the MY82 dissipation is $24/16.6 = 1.45$ times MYNN's.

Physically, \eqref{eq:stress} balances Rotta return-to-isotropy against distortion. The $q/l$ factors hidden in $3l_1/q$ are the inverse of a relaxation timescale
\begin{equation}
\tau \propto \frac{l}{q},
\end{equation}
on which every moment decays toward its isotropic value ($\delta_{ij}q^2/3$ for the stresses, zero for the fluxes), while the mean strain and buoyancy continuously regenerate anisotropy. The $C_1$ terms are the ``isotropization of production'' correction — a fraction of the production tensor fed back into the pressure--strain model. Buoyancy couples the chain downward: the stresses feel the heat flux, the heat flux feels the temperature variance, the temperature variance closes on the heat flux and the mean gradient. Everything in both variants follows from what one is willing to assume when solving \eqref{eq:stress}--\eqref{eq:tvar}.

\section{The master length scale}
\label{sec:lscale}

\key{3}\kh{The scale $l$ is the closure's only free geometric quantity} — it is the size of the energy-containing eddies, and through \eqref{eq:eps} it sets the dissipation. It is also the one part of the scheme that stays one-dimensional: a single scalar for all thirteen moments, built from the vertical structure of the column alone. The default formulation (\texttt{pbl3d\_l\_opt}=1) is Blackadar's interpolation between the surface and an asymptotic value,
\begin{equation}
l = \frac{\kappa z\, l_0}{\kappa z + l_0},
\label{eq:blackadar}
\end{equation}
with $\kappa$ the von-Kármán constant and $z$ the height above ground: near the wall the eddy size is capped by the distance to the wall, $l \to \kappa z$, and far from it $l \to l_0$. The asymptotic scale is the $q$-weighted mean height of the turbulence (Mellor--Yamada 1974, Eq.~72),
\begin{equation}
l_0 = \alpha\,\frac{\displaystyle\int \tilde q\, z\, dz}{\displaystyle\int \tilde q\, dz},
\qquad \alpha = 0.1,
\qquad \tilde q = \max\big(q - q_{\min},\, 0\big),
\label{eq:l0}
\end{equation}
integrated over the column. The weighting by the turbulent excess $\tilde q$ rather than by $q$ itself matters more than it looks: $q^2$ carries a numerical floor $q^2_{\min}$, and with raw-$q$ weighting every quiescent free-atmosphere level contributes the floor value, making $l_0$ scale with the model-top height instead of with the depth of the turbulent layer ($l_0 \to \alpha H/2$ for a quiet column of depth $H$). MYNN solves the same problem by truncating the integral at the diagnosed PBL top; over complex terrain, where a single PBL height is ill-posed, the excess weighting achieves the same end without needing one.

Two limits then act on \eqref{eq:blackadar}. In stable stratification,
\begin{equation}
l \le 0.53\,\frac{q}{N},
\qquad N^2 = \beta g \pd{\Theta_v}{z} > 0,
\label{eq:deardorff}
\end{equation}
the Deardorff (1980) buoyancy limit, with $N$ the Brunt--Väisälä frequency: an eddy of velocity scale $q$ cannot overturn farther than the distance at which the potential-energy cost of displacement consumes its kinetic energy, $l\,N \sim q$. \kh{This limit is an addition of this fork} (\texttt{pbl3d\_n\_tau\_max}, default 0.53): the published Blackadar option had none, so without \eqref{eq:deardorff} $l$ from \eqref{eq:blackadar} is purely geometric and asymptotes to $l_0$ through arbitrarily strong inversions. In level-2 equilibrium the limit binds only above $Ri \sim 1$, i.e.\ in exactly the nocturnal valley air where $l$ and $q$ have decoupled. MYNN's buoyancy scale is $q/N$ at night (Nakanishi and Niino 2009), 1.9 times longer than $0.53\,q/N$ before its harmonic blend with the surface and turbulent-layer scales shortens it — one measurable reason the two schemes need not agree in stable air. \key{9}Second, a grey-zone taper (Honnert et al.\ 2011, in the MYNN form $P_{\sigma}(2.5\,\Delta x / z_i)$; also an addition of this fork — the published scheme admits no $\Delta x$ dependence, because it treats the cell mean as an ensemble mean): the scale-awareness function $P_{\sigma} \in [0,1]$ multiplies the length scales, shrinking the parameterized eddies as the grid begins to resolve them. \kh{In this campaign's nocturnal cases it is inert} (measured: $P_\sigma = 1$ at $\Delta x = 500$~m over a $\sim$100~m stable layer), and by day it is nearly so — from the formula, $0.99$ for $z_i = 1$~km and $0.92$ for 2~km, not yet read from a run — so the grid is ``mesoscale'' with respect to the stable layer even while grey-zone with respect to the daytime convective one. Whether the daytime split between resolved and subgrid kinetic energy is turbulence or grid-set convection is a separate question that the taper does not answer.

Alternatives sit behind \texttt{pbl3d\_l\_opt}: a MYNN-style harmonic blend of surface, turbulent and buoyancy scales, and the Bougeault--Lacarrère parcel formulation, in which the upward and downward displacements a parcel of energy $E(z)$ can afford against the stratification give the mixing scale $\min(l_\mathrm{up}, l_\mathrm{down})$ and the dissipation scale $\sqrt{l_\mathrm{up} l_\mathrm{down}}$.

One rule ties the section together: since the ``one master length scale'' revision in this fork, \emph{whatever limiting is applied to $l$ is seen by both the flux closure and the dissipation \eqref{eq:eps}}. Production and dissipation always refer to the same eddy. The alternative — limiting the scale inside the stress closure while dissipating on the unlimited one — under-pays the dissipation by exactly the limiter's ratio, and a budget that runs a few percent hot for minutes is an exponential.

\section{The approximate closure (\texttt{pbl3d\_opt}$=\pm1$)}
\label{sec:approx}

\subsection{The boundary-layer approximation, applied inside the algebra only}

\key{1}The approximate variant makes one further assumption, and makes it \emph{only inside the algebraic system} \eqref{eq:stress}--\eqref{eq:tvar}: \kh{horizontal gradients of the mean fields, and $\partial W/\partial z$, are dropped there.} This is Mellor's boundary-layer approximation — the same closure content as a 1D level-2.5 scheme, and exactly the simplification whose consequences the campaign wants to isolate. Its payoff is that the system decouples and can be inverted on paper. The two vertical fluxes collapse to eddy-diffusivity form,
\begin{equation}
\avg{uw} = -l\,q\,S_M \pd{U}{z},
\qquad
\avg{vw} = -l\,q\,S_M \pd{V}{z},
\qquad
\avg{w\theta_v} = -l\,q\,S_H \pd{\Theta_v}{z},
\label{eq:kgrad}
\end{equation}
so that $K_m = l q S_M$ and $K_h = l q S_H$ are ordinary eddy diffusivities (m$^2\,$s$^{-1}$), with all of the stability physics packed into the dimensionless functions $S_M$ and $S_H$.

\subsection{The stability functions}

At level 2, where $q$ itself is in equilibrium, $S_M$ and $S_H$ depend only on the gradient Richardson number
\begin{equation}
Ri = \frac{\beta g\, \partial\Theta_v/\partial z}{(\partial U/\partial z)^2 + (\partial V/\partial z)^2},
\end{equation}
through the flux Richardson number $R_f = -P_b/P_s$, obtained from $Ri$ by solving the equilibrium quadratic (Yamada 1975):
\begin{equation}
R_f = c_1\left( Ri + c_2 - \sqrt{Ri^2 + c_3\, Ri + c_4} \right),
\end{equation}
with $c_1 \ldots c_4$ fixed combinations of $(A_1,B_1,A_2,B_2,C_1)$. The functions themselves have the form
\begin{equation}
S_H \propto \frac{R_{f,c}-R_f}{1-R_f},
\qquad
S_M = \frac{A_1}{A_2}\cdot\frac{R_{f,1}-R_f}{R_{f,2}-R_f}\; S_H,
\label{eq:smsh}
\end{equation}
where $R_{f,1}, R_{f,2}$ and the critical value $R_{f,c}$ again follow from the constants; for MY82, $R_{f,c} \approx 0.195$. \key{4}\kh{At $R_f = R_{f,c}$ both functions reach zero and turbulence is cut off entirely} — the closure's own statement that beyond a critical stability (a gradient Richardson number of about 0.2), shear production cannot balance the buoyant destruction plus dissipation. The MYNN control as run has no such cutoff — WRF's version replaces $A_2$ by $A_2/(1+Ri)$ and keeps $S_M, S_H > 0$ at every $Ri$ — and the full closure has none either (Section~\ref{sec:full}: no stability functions, $N$ enters through $l$ only). In the nocturnal valley, where 69~\% of the cells sit above $Ri = 0.3$, the approximate variant is therefore the only one of the three that switches itself off. At level 2.5 the same functions are corrected for disequilibrium following Helfand and Labraga (1988): where the actual $q$ falls below its level-2 equilibrium value $q_\mathrm{eq}$, the functions are rescaled by $q/q_\mathrm{eq}$, which keeps them regular (the naive level-2.5 forms are singular for growing turbulence) at the cost of hard caps on $S_M$ and $S_H$ where the rescaling would overshoot.

\subsection{The other twelve moments}

What distinguishes this variant from MYNN is what happens next. The remaining moments — which a 1D scheme never forms, because homogeneity would erase their divergence anyway — are diagnosed as closed-form by-products of \eqref{eq:kgrad}. The horizontal--horizontal stress and the horizontal heat fluxes are
\begin{equation}
\avg{uv} = \frac{3 A_1 l}{q}\Big( -\avg{uw}\pd{V}{z} - \avg{vw}\pd{U}{z} \Big),
\qquad
\avg{u\theta_v} = 3 A_2\, l^2 \big(S_H + S_M\big) \pd{\Theta_v}{z}\pd{U}{z},
\label{eq:uv}
\end{equation}
(similarly $\avg{v\theta_v}$), the temperature variance is
\begin{equation}
\avg{\theta_v^2} = B_2\, l^2 S_H \left(\pd{\Theta_v}{z}\right)^{\!2},
\end{equation}
and the normal stresses are isotropy plus a strain correction, e.g.
\begin{equation}
\avg{w^2} = \frac{q^2}{3} + \frac{2 A_1 l}{q}\Big( \avg{uw}\pd{U}{z} + \avg{vw}\pd{V}{z} + 2\beta g \avg{w\theta_v} \Big),
\label{eq:w2}
\end{equation}
with analogous expressions for $\avg{u^2}$, $\avg{v^2}$ and the moisture fluxes — fourteen expressions in total. Note what \eqref{eq:w2} says about the physics available: each normal stress is $q^2/3$ plus a perturbation. In stable shear the correction drains $\avg{w^2}$ in favour of the streamwise variance, which is qualitatively right, but \key{6}\kh{anisotropy is a perturbative decoration of an isotropic state, never a solved-for property — the tensor cannot rotate its principal axes into a slope flow.}

The full three-dimensional divergence of all these fluxes is then applied to \eqref{eq:mean}. At \texttt{pbl3d\_opt}$=-1$ only the vertical part is applied and stock Smagorinsky diffusion is retained horizontally, so the $\pm1$ pair separates two questions usually confounded: \emph{does the closure matter, or does applying it in three dimensions matter?} Only the second differs between them.

\subsection{Why this variant is robust}

The approximation buys robustness along with cost. Substituting \eqref{eq:kgrad} into \eqref{eq:ps},
\begin{equation}
P_s = l\,q\,S_M \left[ \left(\pd{U}{z}\right)^{\!2} + \left(\pd{V}{z}\right)^{\!2} \right] \;\ge\; 0,
\label{eq:psapprox}
\end{equation}
a positive-semi-definite quadratic form in the shear: production is manifestly a source, bounded by the capped $S_M$, and the flux is anti-parallel to the gradient by construction. Together with the $R_{f,c}$ cutoff this makes the variant hard to destabilize — the same structural guarantees that two decades of operational 1D Mellor--Yamada schemes rest on. It runs at roughly the cost of MYNN (about 0.8~s per step on this grid) and serves as the campaign's \emph{diagnostic}: it shares the coordinate, the gradients, the length scales and the plumbing of the full closure while retaining the analytic solvability of a 1D scheme, so a failure present in one branch and absent in the other is localized to the flux solve itself.

\section{The full closure (\texttt{pbl3d\_opt}$=2$)}
\label{sec:full}

\subsection{The linear system}

The full variant makes no homogeneity assumption anywhere. All nine mean-velocity gradients and all three gradients of $\Theta_v$ (and $Q_v$) stay in \eqref{eq:stress}--\eqref{eq:tvar}; the system no longer decouples and there is no closed form. It is solved numerically at every grid point where $q^2$ exceeds its floor, at every substep: a dense $10\times10$ linear system
\begin{equation}
\mathsf{A}\,\mathbf{x} = \mathbf{b},
\qquad
\mathbf{x} = \big( \avg{u^2}, \avg{v^2}, \avg{w^2},\; \avg{uv}, \avg{uw}, \avg{vw},\;
\avg{u\theta_v}, \avg{v\theta_v}, \avg{w\theta_v},\; \avg{\theta_v^2} \big)^{\!\top},
\end{equation}
followed by a $4\times4$ system for $\big(\avg{u q_v}, \avg{v q_v}, \avg{w q_v}, \avg{q_v\theta_v}\big)$ — the joint system is block-triangular, so the moisture block can be solved after the momentum--heat block, and the moisture variance follows in closed form,
\begin{equation}
\avg{q_v^2} = -\frac{B_2 l}{q}\, \avg{u_j q_v} \pd{Q_v}{x_j}.
\end{equation}

The matrix has the structure $\mathsf{A} = \tau^{-1}\mathsf{I} + \mathsf{S}$: return-to-isotropy rates on the diagonal, mean strain off it, buoyancy coupling the blocks. Representative entries make the anatomy concrete:
\begin{equation}
\mathsf{A}_{11} = \frac{q}{2A_1 l} + 2\pd{U}{x},
\qquad
\mathsf{A}_{55} = \frac{q}{3A_1 l} + \pd{W}{z} + \pd{U}{x},
\qquad
\mathsf{A}_{99} = \frac{q}{3A_2 l} + \pd{W}{z},
\qquad
\mathsf{A}_{10,10} = \frac{q}{B_2 l},
\end{equation}
\begin{equation}
\mathsf{A}_{1,9} = \beta g,
\qquad
\mathsf{A}_{3,9} = -2\beta g,
\qquad
\mathsf{A}_{9,10} = -\beta g .
\end{equation}
Row 3 ($\avg{w^2}$) is coupled twice as strongly to the heat flux as row 1, and with opposite sign — buoyancy moves energy between the vertical and horizontal variances, it does not create $q^2$ by itself in this system. Only the first six entries of the right-hand side are non-zero: the isotropic energy supply plus the isotropization-of-production terms,
\begin{equation}
b_{1} = \frac{q^3}{6 A_1 l} + 3 C_1 q^2 \pd{U}{x},
\qquad
b_{4} = C_1 q^2 \left( \pd{V}{x} + \pd{U}{y} \right),
\qquad
b_{7} = \dots = b_{10} = 0,
\end{equation}
with $b_2, b_3$ and $b_5, b_6$ the obvious permutations. The heat-flux and variance rows are driven entirely through the coupling — they have no source of their own, which is the algebraic statement that scalar fluxes exist only because stresses act on scalar gradients.

Two properties of $\mathsf{A}$ dictate the numerics. Its columns carry mixed units (s$^{-1}$ on the momentum diagonal, buoyancy entries in m\,s$^{-2}$\,K$^{-1}$), so the raw condition number is meaningless; and near the turbulence cutoff the diagonal shrinks with $q$ while the strain does not, so the system can approach singularity in physically strained, weakly turbulent air. The solve is therefore LAPACK \texttt{dgesvx} in double precision inside a single-precision model, with row/column equilibration requested and the returned condition estimate checked against a ceiling.

\subsection{Realizability, and the safeguard stack}

Because the solved tensor is genuinely anisotropic, nothing in exact arithmetic — let alone finite arithmetic on a near-singular system — guarantees that it is a physically possible one: a variance can come back negative, a correlation can exceed one. And nothing bounds the production such a tensor feeds back into \eqref{eq:qsq}. A stack of safeguards therefore stands between the solver and the model, escalating from prevention to repair.

\emph{First, a strain limiter} in the spirit of Durbin (1996). With the strain-rate magnitude
\begin{equation}
|S| = \sqrt{2\, S_{ij} S_{ij}}
    = \left[ 2\!\left( U_x^2 + V_y^2 + W_z^2 \right)
      + \left( U_y{+}V_x \right)^2 + \left( U_z{+}W_x \right)^2 + \left( V_z{+}W_y \right)^2 \right]^{1/2},
\end{equation}
the scale entering the solve is capped so that the non-dimensional distortion cannot exceed a ceiling $\Sigma_{\max}$ (\texttt{pbl3d\_sk\_eps\_max}, default 6):
\begin{equation}
\frac{|S|\,k}{\epsilon} = \frac{|S|\, B_1 l}{2 q} \le \Sigma_{\max}
\quad\Longleftrightarrow\quad
l_\mathrm{use} = \min\!\left( l,\ \frac{2\,\Sigma_{\max}}{B_1}\,\frac{q}{|S|} \right).
\label{eq:tier1}
\end{equation}
\key{5}Physically: when the mean flow distorts an eddy much faster than the eddy can relax ($|S|\tau \gg 1$), \kh{the algebraic equilibrium assumption behind \eqref{eq:stress} has no basis}, and the solved tensor leaves the realizable region with an unphysically large implied eddy viscosity. Capping $l$ shortens $\tau$ until distortion and relaxation are back in ratio. Because $l_\mathrm{use}$ enters both $\mathsf{A}$ and $\mathbf{b}$, the strain-free isotropic limit is preserved exactly, and — the one-scale rule of Section~\ref{sec:lscale} — the same limited scale reaches the dissipation, so a limited cell transports \emph{and} dissipates at the reduced eddy size. This coupling is load-bearing: with the limiter acting on production only, a cell under strong strain runs its budget a steady 10--20\% hot, which is a $q^2$ exponential with a doubling time of about a minute; with the consistent scale the same cell's production-to-dissipation ratio sits below one and the runaway never starts.

\emph{Second, back off and re-solve.} If the solve is degenerate ($\texttt{info} \neq 0$, condition estimate above the ceiling) or the returned tensor is not positive semi-definite, $l$ is halved and the system solved again, up to six times. Each retry is an \emph{exact} solution of a well-posed problem at a shorter eddy scale — not an approximate solution of the original — which is what makes shrinking $l$ an admissible response rather than a fudge.

\emph{Third, project into the realizable set} whatever still emerges flawed: floor the variances; rescale to the exact trace the system implies (summing the three variance rows of $\mathsf{A}\mathbf{x}=\mathbf{b}$ yields an identity the solution must satisfy); enforce the Cauchy--Schwarz bounds
\begin{equation}
\avg{u_i u_j}^2 \le \avg{u_i^2}\avg{u_j^2},
\qquad
\avg{u_i \theta_v}^2 \le \avg{u_i^2}\avg{\theta_v^2},
\qquad
\avg{u_i q_v}^2 \le \avg{u_i^2}\avg{q_v^2},
\end{equation}
and complete positive semi-definiteness with a determinant-based shrink of the off-diagonal block. The same constraints cover the moisture moments — an implied correlation between $q_v$ and $\theta_v$ fluctuations cannot exceed one any more than a momentum correlation can.

\subsection{What the full production changes}

The $q^2$ budget \eqref{eq:qsq} gains the complete nine-term contraction
\begin{equation}
P_s = -\avg{u^2}\pd{U}{x} - \avg{v^2}\pd{V}{y} - \avg{w^2}\pd{W}{z}
      - \avg{uv}\!\left(\pd{U}{y}{+}\pd{V}{x}\right)
      - \avg{uw}\!\left(\pd{U}{z}{+}\pd{W}{x}\right)
      - \avg{vw}\!\left(\pd{V}{z}{+}\pd{W}{y}\right),
\label{eq:ps3d}
\end{equation}
which on a $34^\circ$ slope is dominated by cross terms a 1D scheme never sees. \key{6}Unlike the approximate variant's quadratic form \eqref{eq:psapprox}, \eqref{eq:ps3d} has no built-in sign: \kh{with a fully anisotropic tensor, locally negative production} — energy returned from the turbulence to the mean flow — is physically admissible and does occur, most commonly where the stress carried over from the previous state opposes a rotated strain. The scheme must therefore tolerate a production term that is a genuine two-way exchange, and the level-2 diagnostic branch guards its $q^3 = B_1 l\,(P_s + P_b)$ inversion with an explicit $\max(\,\cdot\,,0)$ before taking the two-thirds power. Horizontal diffusion of $q^2$ itself, mirroring the resolved-scale operator with its own halo exchange, is also active only on this path.

The measured cost on this grid is 1.3--1.4~s per step against $\sim$0.8 for the approximate variant — the dense solves plus a fraction of a repeat solve per point from the back-off ladder — consistent with the 20--50\% over MYNN quoted by Juliano et al.

\section{Numerics shared by both variants}
\label{sec:numerics}

Downstream of the closure the two variants are identical, and three implementation choices shape how either behaves in practice.

\emph{The vertical flux divergence is explicit.} Where a 1D PBL scheme hides its stiffness in an implicit tridiagonal solve per column, this scheme applies
\begin{equation}
\left(\pd{U}{t}\right)_{\!\mathrm{turb}} = -\frac{1}{\rho}\,\pd{\big(\rho \avg{uw}\big)}{z}
\end{equation}
(and its analogues) as an explicit update, sub-stepped \texttt{pbl3d\_nsteps} times per Runge--Kutta step. Explicit diffusion is conditionally stable — per substep, effectively $K\,\Delta t/\Delta z^2 \lesssim \tfrac12$ — so with $\Delta z \approx 17$~m near the ground, strong mixing episodes are absorbed by the substepping rather than by damping — the scheme meets a violent flux with more, smaller steps, not with the unconditional stability of an implicit solve.

\emph{The horizontal derivatives are metric-corrected.} In the terrain-following coordinate $\eta$, a constant-height derivative is assembled as
\begin{equation}
\left.\pd{U}{x}\right|_{z}
 = m \left.\pd{(U/m)}{x}\right|_{\eta}
 - \left.\pd{z}{x}\right|_{\eta}\, \pd{U}{z},
\end{equation}
with $m$ the map-scale factor — on a $30^\circ$ slope the correction term is comparable to the raw derivative, and omitting it would manufacture strain from the coordinate itself. The horizontal flux divergences carry the same correction, and are damped on steep slopes by a factor
\begin{equation}
\sigma = \max\!\left( \sqrt{ z_x^2 + z_y^2 },\ 1 \right)
\end{equation}
dividing the horizontal tendencies, a pragmatic guard where the coordinate surfaces fold too steeply for the discretization. (In this fork the energy paid into $q^2$ by the horizontal shear production is paired with the slope-tapered tendency the flow actually feels, so the taper cannot become a spurious energy source.)

\key{8}\emph{The surface supplies a stress, not a $q^2$.} \kh{The lower boundary is Monin--Obukhov similarity on a flat surface}: the surface-layer scheme hands over the friction velocity $u_*$ (m\,s$^{-1}$), and the ground face of the stress is set to $\avg{uw} = -u_*^2$, rotated into the wind. Nothing anchors $q^2$ there. The neutral wall equilibrium — production balancing dissipation at the lowest level, $u_*^3/(\kappa z) = q^3/(B_1 \kappa z)$, hence
\begin{equation}
q^2 \ge B_1^{2/3}\, u_*^2 \approx 6.5\,u_*^2
\label{eq:wall}
\end{equation}
— is imposed by every MYNN-family scheme and is absent from the inherited closure. (\emph{Correction:} the 2026-09-01 version of this note stated \eqref{eq:wall} as active. It exists in this fork as an opt-in switch, \texttt{pbl3d\_sfc\_qsq\_bc}, default off, tested once so far: neutral in the windy evening transition of 18 July 2025; a retest on a calm night is pending.) Two further facts belong to the configuration as run: the surface-layer scheme floors $u_*$ at 0.03~m\,s$^{-1}$, which is active in 18~\% of land cells at night with no measured footprint in $q^2$; and the three-dimensional diagnostic surface layer of Eghdami et al.\ (2022), which would replace the flat-surface similarity at the bottom of the algebraic system, exists as an option (\texttt{pbl3d\_sfc\_opt}) and is off.

\section{Side by side}

\begin{center}
\footnotesize
\begin{tabular}{@{}lll@{}}
\toprule
 & Approximate 3D (\texttt{pbl3d\_opt}$=\pm1$) & Full 3D (\texttt{pbl3d\_opt}$=2$) \\
\midrule
Closure content & Level 2.5 under horizontal homogeneity & Level 2.5, no homogeneity assumption \\
Gradients in the closure & 2 of 9 (via $S_M$, $S_H$) & 9 of 9 (12 with the scalar gradients) \\
Second moments & 14 closed forms; anisotropy diagnosed & $10{\times}10$ + $4{\times}4$ solve; anisotropy solved \\
Normal stresses & $q^2/3$ + perturbation \eqref{eq:w2} & Fully coupled unknowns \\
$q^2$ shear production & Quadratic form \eqref{eq:psapprox}, $\ge 0$ & Nine-term contraction \eqref{eq:ps3d}, unbounded sign \\
Stability guard & $R_{f,c}$ cutoff + Helfand--Labraga caps & Strain limiter \eqref{eq:tier1}, re-solve, projection \\
Turbulence cutoff & Yes: $S_M = S_H = 0$ at $R_f \ge 0.195$ ($Ri \approx 0.2$) & None (nor in the MYNN control as run) \\
Flux divergence applied & 3D at $+1$ (vertical only at $-1$) & 3D always \\
Cost (this grid) & $\sim$0.8 s/step, $\approx$ MYNN & 1.3--1.4 s/step \\
Role in the campaign & The diagnostic control & The object of study \\
\bottomrule
\end{tabular}
\end{center}

The essential picture: there is \emph{one} closure — Mellor--Yamada level 2.5 with prognostic $q^2$ and one master length scale — and one design decision: whether the algebraic second-moment system is simplified until it inverts analytically, or solved as it stands. The approximate variant keeps the guarantees of the 1D tradition while already exercising the three-dimensional \emph{application} of the fluxes, which makes it the clean control for every question about coordinates, divergences and slopes. The full variant is the only one whose stress tensor can tilt its principal axes into a katabatic jet and pay energy back to the mean flow where the physics says it should — and the price of removing the homogeneity assumption is that boundedness and realizability, which the analytic solution provides for free, must be engineered around the solver.


\section{The assumptions the PhD concept rests on}
\label{sec:concept}

The concept draft (RQ~I: the MYNN control, the Goger-2019 port, the approximate and the full 3D closure at $\Delta x = 500$~m, judged by regime against the TEAMx observations; RQ~II: a large-eddy simulation as reference for the UAS turbulence) is read here against the assumptions of Sections~\ref{sec:core}--\ref{sec:numerics}. An assumption carries a \keytag{n} in the margin above wherever a research question, a regime split or a decision metric of the concept turns on it. Table~\ref{tab:keys} collects the nine; the two subsections after it say which of the concept's choices are sound and what is missing. M, G, A, F are the four configurations of the concept; the last column points to the row of the companion note \emph{The assumptions of a one-dimensional PBL closure, and which scheme drops which} (2026-09-07), where the equations and code lines are.

\begin{table}[!t]
\centering
\caption{The nine assumptions the concept turns on. K = kept, P = relaxed or partly dropped, D = dropped, for the MYNN control as run (M), the Goger-2019 port (G), the approximate (A) and the full (F) 3D closure. ``Measured'' means read from this project's runs or from the run's namelist; ``inferred'' means derived from the code or a formula; ``unrun'' means the switch exists and has not been exercised.}
\label{tab:keys}
\footnotesize
\setlength{\tabcolsep}{4pt}
\renewcommand{\arraystretch}{1.12}
\begin{tabularx}{\textwidth}{@{}l >{\raggedright\arraybackslash}p{3.9cm} >{\raggedright\arraybackslash}X c@{\,}c@{\,}c@{\,}c >{\raggedright\arraybackslash}X l@{}}
\toprule
 & Assumption & Where the concept turns on it & M & G & A & F & Evidence in this project & Row \\
\midrule
\keytag{1} & Horizontal gradients of the mean fields, and the gradients of $W$, produce no turbulence and do not enter the flux relations. & The nine-term shear-production matrix that organises the state of research; the M--G--A--F ladder; the shear partition of RQ~I. & K & P & K & D & After the pairing fix the six horizontal terms carry $\sim$5--10~\% of the production (inferred); the nine shares by regime have never been tabulated. & B1, B2 \\
\addlinespace[2pt]
\keytag{2} & The mean flow feels turbulence only through the vertical flux divergence; horizontal exchange is a numerical operator with the grid as its length. & The control keeps Smagorinsky so that the ladder measures the closure; the question of Xu et al.\ (2026) — length scale or horizontal flux at night. & K & K & D & D & The horizontal flux divergence against the resolved advection, binned by $Ri$, is not yet measured (WRFlux stores both). & B3, B5 \\
\addlinespace[2pt]
\keytag{3} & One master length scale for transport and dissipation, built from the vertical structure alone (Blackadar with the excess-weighted $l_0$, the buoyancy limit, the strain cap). & ``Length scales evaluated as part of each scheme''; ``the last 1D relic''; the stable-regime question ``length scale or constants''. & K & P & K & K & The nocturnal deficit (0.16 of MYNN) is an equilibrium, flat above $Ri = 0.3$ (measured); which bound sets $l$ at night is inferred from the code, not read from the arrays; the MYNN-blend option crashed the explicit vertical diffusion after 14~min and is unrun. & C4, C5 \\
\addlinespace[2pt]
\keytag{4} & Turbulence dies at a critical Richardson number. & Not mentioned in the concept. & D & P & K & D & 69~\% of the nocturnal valley cells above $Ri = 0.3$ (measured, full closure); the fraction of cells where A has $S_M = 0$ is not measured. & C6 \\
\addlinespace[2pt]
\keytag{5} & All moments but $q^2$ are in local equilibrium with the mean gradients; the relaxation time $l/q$ is short against the flow's changes. & The evening transition as a regime; the ``limiter footprint'' metric. & K & K & K & K & The strain cap is load-bearing: $\Sigma_{\max} = 12$ or off restores the nocturnal runaway within 45~min (measured). & C7 \\
\addlinespace[2pt]
\keytag{6} & The flux is aligned with the gradient of its own quantity (K-form); anisotropy is at most a perturbation of isotropy. & ``$K_M$ from observed first and second moments vs modelled $K_M$''; the katabatic-jet argument for the full closure. & P & K & P & D & The solved tensor's anisotropy has never been compared with the i-Box sonics. & C1, C3 \\
\addlinespace[2pt]
\keytag{7} & The closure constants are universal. & ``Length scale or constants''; the rule not to tune before the observational hold. & K & K & K & K & Only MY82 has been run; $B_1$ alone makes the two families dissipate 1.45 apart at equal $q$ and $l$ (from the constants). & C9 \\
\addlinespace[2pt]
\keytag{8} & The lower boundary is a flat, homogeneous surface: Monin--Obukhov fluxes, one $u_*$ and one roughness per cell, no wall value of $q^2$. & Methods: ``3D diagnostic surface layer as the scheme's bottom boundary''; the $u_*$ floor. & K & K & K & K & Sonics hold $q^2 = 0.13$--$0.41$~m$^2$\,s$^{-2}$ all night at seven i-Box sites while every model goes laminar (measured); the 3D surface option is unrun; the wall condition \eqref{eq:wall} is neutral in the windy transition (measured). & B6 \\
\addlinespace[2pt]
\keytag{9} & The cell mean is an ensemble mean and every energetic eddy is subgrid; the closure knows no $\Delta x$. & Convective regime judged on subgrid plus resolved TKE; the effective resolution $\sim 7\,\Delta x$; the partition scaling behind RQ~II. & P & K & P & P & Taper inert at night (measured), $\le 8$~\% on $l$ by day (formula); 85--91~\% of the kinetic energy resolved at 43~m in the morning mixed layer, unproven as turbulence — the $w$-spectrum test is pending. & A1, A2 \\
\bottomrule
\end{tabularx}
\end{table}

\subsection{What the concept gets right}

\emph{The nine-term matrix as the organising object} (\keytag{1}). It is the one assumption whose removal is the full closure's reason to exist, and the one the HOTSPOT observations can reach — four UAS positions across the valley and the lidar RHI scans give the horizontal gradients that no tower network could. Sound, and the right red thread for a shear-financed project.

\emph{Keeping Smagorinsky in the control} (\keytag{2}). Without it the step from the control to the 3D variants would change the horizontal exchange twice over; with it the ladder measures the closure. The plan to hold each closure's horizontal flux divergence against the resolved advection in the WRFlux budget, binned by $Ri$, is the test Xu et al.\ (2026) could not run.

\emph{Length scales evaluated per scheme, nothing tuned before the observational hold} (\keytag{3}, \keytag{7}). This matches the literature — at 750~m in a stable cold pool the scale, not the horizontal flux, was the lever (Arthur et al.\ 2022) — and this project's own finding that the nocturnal deficit is an equilibrium, not spin-up.

\emph{Regimes by Richardson number and flow, not by the sign of the surface heat flux; subgrid plus resolved TKE in unstable air; momentum flux as the sharper discriminator; $K_M$ from observed first and second moments} (\keytag{6}, \keytag{9}). The last is the one direct test of the K-form assumption, and only the UAS with the budget station can deliver it. Measured and inferred are kept apart in the preliminary findings.

\subsection{What is missing or needs sharpening}

\begin{enumerate}\setlength{\itemsep}{2pt}
\item \textbf{The cutoff nobody mentions} (\keytag{4}). The approximate variant carries the level-2 cutoff, $S_M = S_H = 0$ at $R_f \ge 0.195$; the control as run has none ($A_2 \to A_2/(1+Ri)$), the full closure has none. In a nocturnal valley with 69~\% of the cells above $Ri = 0.3$, the approximate variant's night is a cutoff experiment before it is a closure experiment, and its stable-regime result cannot be read as ``the boundary-layer approximation'' alone. Table~1 of the concept needs a ``stability cutoff'' column, RQ~I(iii) one sentence, and the paper one number: the fraction of nocturnal valley cells in which A has switched itself off, against the fraction in which F is realizable.

\item \textbf{The control as run is not the control the methods describe} (\keytag{6}, \keytag{9}). Measured from the namelist of job 8320565: no \texttt{bl\_mynn} key is set, so WRF~4.8's defaults apply — closure level 2.6, mass-flux transport of heat, moisture \emph{and momentum} on, TKE advection off. The methods bullet claims momentum mass flux off and TKE advection on. Either the control is rerun or the text is corrected; the companion note has it right. Beyond the wording: the mass-flux plume is a nonlocal transport that none of the three other schemes has, so in the convective regime the step M$\to$F changes the transport model and the dimensionality at once. A second control — mass flux off, TKE advection on, level 2.5 — costs one run and turns M$\to$A into a one-assumption step; it is the cheapest experiment in the concept and is not in it.

\item \textbf{Name the bound that sets $l$ at night, and factorise the attribution} (\keytag{3}). The concept's ``$l_0$ too small'' is Juliano's convective statement. At night the fork's $l$ is set by the buoyancy limit $0.53\,q/N$ and the strain cap, both fork additions (Section~\ref{sec:lscale}) that the paper must describe as such. The attribution ``length scale or constants'' does not have to stay a hypothesis: in the level-2 balance the control and the approximate variant share,
\begin{equation}
q^2_{\mathrm{eq}} = B_1\, S_M(Ri)\,(1 - R_f)\; l^2 \left[\left(\pd{U}{z}\right)^{\!2} + \left(\pd{V}{z}\right)^{\!2}\right],
\label{eq:factor}
\end{equation}
so at equal shear the ratio of two schemes' equilibrium $q^2$ is the product of a constants factor ($B_1$: $16.6/24 = 0.69$), a stability-function factor and the squared length-scale ratio — three stored fields, hence a per-$Ri$-bin measurement (for the full closure with $S_M$ read back from the solved stress). The concept should name \eqref{eq:factor} as the paper-1 attribution method. It should also carry the risk the code has already shown: the longer MYNN-blend scale drove the explicit vertical diffusion past its stability limit (diffusion number $K\Delta t/\Delta z^2 > 0.5$ in 2.4~\% of the lowest levels, crash after 14~min), so the length-scale experiment is gated by numerics — an implicit vertical solve or a diffusion-number cap is a prerequisite that belongs in the schedule, not a detail.

\item \textbf{The surface is described as it is not run} (\keytag{8}). The methods bullet reads as if the Eghdami et al.\ (2022) surface layer is the bottom boundary; it is off. The configuration as run is flat-surface similarity with a $u_*$ floor of 0.03~m\,s$^{-1}$ and no wall value of $q^2$ — the one condition every MYNN-family scheme imposes and the inherited closure lacks (Section~\ref{sec:numerics}). The surface layer also deserves promotion from a methods detail to a hypothesis of regime (iii): the largest closure-independent discrepancy of the project is measured there — sonic $q^2$ of 0.13--0.41~m$^2$\,s$^{-2}$ all night at seven sites while every model, MYNN included, goes laminar — and it is where the literature's warnings (a constant-flux layer in only 8--28~\% of the daytime periods; a surface layer stable 4.7~\% of the time under an often stable column at Innsbruck) meet the project's own 0--100~m problem. RQ~I(iii) should say that the nocturnal $q^2$ is judged against the sonics — the observed $u_*$ of 0.07--0.16~m\,s$^{-1}$ implies a wall $q^2$ of 0.03--0.17, between the two closures — and not against the control.

\item \textbf{The grey-zone premise, and what the convective result has not yet shown} (\keytag{9}). The published scheme is grid-free because it declares the cell mean an ensemble mean; the taper and the buoyancy limit are fork additions, so ``the closure knows no $\Delta x$'' is true of the scheme as published and nearly true of the scheme as run (taper $\ge 0.92$ by day, from the formula). The concept's claim that the subgrid ratio of 0.28 in the morning mixed layer is ``a partition, not a deficit'' rests on 85--91~\% of the kinetic energy being resolved at 43~m — the signature of grid-set convection as much as of turbulence, since coarse-grained LES expects about three quarters of mixed-layer TKE subgrid at $\Delta x/(z_i + h_c) \approx 0.6$. The analysis section lists the effective-resolution check; the claim should be tied to it: a $w$-spectrum with energy below $6$--$8\,\Delta x$ is the condition under which the partition reading stands.

\item \textbf{Three constant sets in one comparison} (\keytag{7}). NN09 in M and G, MY82 in A and F, the Bougeault--Lacarrère-tuned set in the publication. The paper must state that the schemes differ in constants — the $B_1$ factor of 1.45 in the dissipation at equal $q$ and $l$ is larger than several of the effects it will report — before it attributes a difference to the tensor.

\item \textbf{Assumptions the concept does not name.} (a)~Local equilibrium (\keytag{5}): the evening transition is exactly the regime in which the relaxation time $l/q$ (minutes) is not short against the flow (an hour); the concept lists the transition and the limiter footprint but does not connect them — regime (ii) tests \keytag{5}, and the cap footprint by regime is its measurement. (b)~Pressure transport absorbed into down-gradient transport, and Rotta's return to isotropy with the isotropisation constant $C_1$: kept by every scheme in the comparison and testable by one observation only, the residual of the Radfeld budget station; the concept should say so where it lists the budget terms. (c)~Buoyancy in the stress equations — the three heat-flux components of the full closure, the slope-normal surface heat flux with its horizontal component — has no home in the concept (its own open item~9); one bullet in the 3D-closure subsection would close it.

\item \textbf{Fork versus publication.} The methods say the 3D closure switches ``all default to the reference behaviour''. That is true of the Registry and false of the production namelist: the runs use one master scale, the strain cap, the buoyancy limit, the excess-weighted $l_0$, the slope-factor energy pairing and the $u_*$ floor — six departures from Juliano et al.\ (2022), each recorded with its reason, none of them in the concept's Table~1. Paper~1 compares the scheme as modified; the table's F row needs the list, and each scheme an ``as run'' note.
\end{enumerate}

Net: the assumption set behind RQ~I is the right one — the rows that decide the questions, \keytag{1}, \keytag{2}, \keytag{3} and \keytag{6}, are all in the concept, and the tests it names for them are the right tests. What the draft describes, though, are the published schemes; the schemes that will run carry six modifications, one confounded control and one cutoff nobody has mentioned. Correcting the text is an afternoon; the second control is one run.

\begin{multicols}{2}
\begin{thebibliography}{99}\footnotesize\setlength{\itemsep}{0pt}\setlength{\parskip}{0pt}
\bibitem{MY74} Mellor, G.~L., and T.~Yamada, 1974: A hierarchy of turbulence closure models for planetary boundary layers. \emph{J.~Atmos.~Sci.}, \textbf{31}, 1791--1806.
\bibitem{Y75} Yamada, T., 1975: The critical Richardson number and the ratio of the eddy transport coefficients obtained from a turbulence closure model. \emph{J.~Atmos.~Sci.}, \textbf{32}, 926--933.
\bibitem{Deardorff80} Deardorff, J.~W., 1980: Stratocumulus-capped mixed layers derived from a three-dimensional model. \emph{Bound.-Layer Meteor.}, \textbf{18}, 495--527.
\bibitem{MY82} Mellor, G.~L., and T.~Yamada, 1982: Development of a turbulence closure model for geophysical fluid problems. \emph{Rev.~Geophys.~Space Phys.}, \textbf{20}, 851--875.
\bibitem{HL88} Helfand, H.~M., and J.~C.~Labraga, 1988: Design of a nonsingular level 2.5 second-order closure model for the prediction of atmospheric turbulence. \emph{J.~Atmos.~Sci.}, \textbf{45}, 113--132.
\bibitem{BL89} Bougeault, P., and P.~Lacarrère, 1989: Parameterization of orography-induced turbulence in a mesobeta-scale model. \emph{Mon.~Wea.~Rev.}, \textbf{117}, 1872--1890.
\bibitem{Durbin96} Durbin, P.~A., 1996: On the $k$--$\varepsilon$ stagnation point anomaly. \emph{Int.~J.~Heat Fluid Flow}, \textbf{17}, 89--90.
\bibitem{NN09} Nakanishi, M., and H.~Niino, 2009: Development of an improved turbulence closure model for the atmospheric boundary layer. \emph{J.~Meteor.~Soc.~Japan}, \textbf{87}, 895--912.
\bibitem{Honnert11} Honnert, R., V.~Masson, and F.~Couvreux, 2011: A diagnostic for evaluating the representation of turbulence in atmospheric models at the kilometric scale. \emph{J.~Atmos.~Sci.}, \textbf{68}, 3112--3131.
\bibitem{Kosovic20} Kosović, B., P.~Jiménez Munoz, T.~W.~Juliano, A.~Martilli, M.~Eghdami, A.~P.~Barros, and S.~E.~Haupt, 2020: Three-dimensional planetary boundary layer parameterization for high-resolution mesoscale simulations. \emph{J.~Phys.:~Conf.~Ser.}, \textbf{1452}, 012080.
\bibitem{Juliano22} Juliano, T.~W., B.~Kosović, P.~A.~Jiménez, M.~Eghdami, S.~E.~Haupt, and A.~Martilli, 2022: ``Gray zone'' simulations using a three-dimensional planetary boundary layer parameterization in the WRF model. \emph{Mon.~Wea.~Rev.}, \textbf{150}, 1585--1619.
\bibitem{Arthur22} Arthur, R.~S., T.~W.~Juliano, B.~Adler, R.~Krishnamurthy, J.~K.~Lundquist, B.~Kosović, and P.~A.~Jiménez, 2022: Improved representation of horizontal variability and turbulence in mesoscale simulations of an extended cold-air pool event. \emph{J.~Appl.~Meteor.~Climatol.}, \textbf{61}, 685--707.
\bibitem{Eghdami22} Eghdami, M., A.~P.~Barros, P.~A.~Jiménez, T.~W.~Juliano, and B.~Kosović, 2022: Diagnosis of second-order turbulent properties of the surface layer for three-dimensional flow based on the Mellor--Yamada model. \emph{Mon.~Wea.~Rev.}, \textbf{150}, 1003--1021.
\bibitem{Xu26} Xu, H., C.~Han, Y.~Zhang, X.~Zhang, T.~W.~Juliano, and Y.~Ma, 2026: Added value of three-dimensional planetary boundary layer schemes for simulating local-scale circulations over complex terrain on the Tibetan Plateau. \emph{J.~Geophys.~Res.~Atmos.}, \textbf{131}, e2025JD046165.
\end{thebibliography}
\end{multicols}

\end{document}
